\section{Test}

Per il test del sistema A.L.I.C.E. sono stati utilizzati i seguenti strumenti:

\begin{itemize}
    \item \textbf{Postman} -- Applicazione utilizzata per effettuare test funzionali sulle API di Supabase, verificandone il corretto funzionamento simulando richieste HTTP (GET, POST, PATCH, DELETE) verso gli endpoint REST auto-generati da PostgREST.

    \item \textbf{ESLint} -- Strumento di analisi statica del codice JavaScript/TypeScript, utilizzato per individuare e correggere bug, errori di sintassi e problemi di stile, migliorando la qualità e la coerenza del codice sorgente.

    \item \textbf{Supabase Dashboard (SQL Editor)} -- Interfaccia web fornita da Supabase per eseguire query SQL dirette sul database PostgreSQL, utilizzata per verificare la corretta persistenza dei dati e il funzionamento delle policy Row Level Security.

    \item \textbf{Chrome DevTools} -- Strumenti di sviluppo integrati nel browser Google Chrome, utilizzati per il debug del frontend React, l'analisi delle richieste di rete, il controllo del rendering dei componenti e la verifica della responsività dell'interfaccia.
\end{itemize}

\newpage
